\documentclass[10pt]{article}
\usepackage[T1]{fontenc}
\usepackage[margin=0.78in]{geometry}
\usepackage{amsmath,amssymb,amsthm}
\usepackage{graphicx}
\usepackage[colorlinks=true,linkcolor=blue,citecolor=blue,urlcolor=blue]{hyperref}
\newtheorem{theorem}{Theorem}
\newtheorem{lemma}[theorem]{Lemma}
\newcommand{\B}{\mathbf B}
\newcommand{\R}{\mathbf R}
\newcommand{\E}{\mathcal E}
\newcommand{\K}{\mathcal K}
\newcommand{\Pcal}{\mathcal P}
\newcommand{\Fock}{\mathcal F}
\title{Harnack and Poisson metrics coincide\\with at most one nonscalar row}
\author{Run \texttt{fa\_banach\_001} --- \texttt{agent\_lane\_00}}
\date{26 August 2026}
\begin{document}
\maketitle

\noindent\textbf{Status.} Likely valid partial result, subject to human
review. It gives an affirmative answer to a dimension class strictly larger
than the noncommutative-ball and regular-polydisk cases stated in the source.

\section{Source question}

Popescu's paper \emph{Hyperbolic geometry on noncommutative polyballs},
arXiv:1701.00766, asks on page 4 whether its Harnack metric $\delta_H$ and
Poisson metric $\delta_{\Pcal}$ coincide on every regular polyball
\cite{metric-source}. The relevant source page is reproduced below; the open
problem is the final paragraph before Section 1.

\begin{center}
\includegraphics[width=.91\textwidth]{figures/open_problem_crop.png}
\end{center}

\noindent\textbf{Transcription.} ``It remains an open problem whether the
same result is true for any regular polyball.'' Here ``the same result'' is
$\delta_H=\delta_{\Pcal}$, known there for a noncommutative ball and for a
regular polydisk.

\section{Definitions and statement}

Let $\B_{\mathbf n}$ be the regular polyball with
$\mathbf n=(n_1,\ldots,n_k)$, and let
$\boldsymbol{\Pcal}(\R,\mathbf X)$ be its free pluriharmonic Poisson kernel.
For $\mathbf A,\mathbf B\in\B_{\mathbf n}(\mathcal H)^-$ and $c>1$, write
$\mathbf A\sim_{\Pcal,c}\mathbf B$ when
\[
 c^{-2}\boldsymbol{\Pcal}(\R,r\mathbf B)
 \le \boldsymbol{\Pcal}(\R,r\mathbf A)
 \le c^2\boldsymbol{\Pcal}(\R,r\mathbf B),\qquad 0\le r<1.
\tag{1}
\]
Write $\mathbf A\sim_{H,c}\mathbf B$ when the same inequalities hold after
replacing the Poisson kernel by every positive operator-valued free
$k$-pluriharmonic function $F$. The metrics are the logarithms of the infima
of the admissible constants on their respective equivalence parts.

\begin{theorem}\label{thm:metric}
Suppose that at most one coordinate of $\mathbf n$ is larger than one. For
every $\mathbf A,\mathbf B\in\B_{\mathbf n}(\mathcal H)^-$ and every $c>1$,
\[
 \mathbf A\sim_{H,c}\mathbf B
 \quad\Longleftrightarrow\quad
 \mathbf A\sim_{\Pcal,c}\mathbf B.
\]
Consequently the Harnack parts and Poisson parts coincide, and
\[
 \delta_H(\mathbf A,\mathbf B)
 =\delta_{\Pcal}(\mathbf A,\mathbf B)
\]
on every such part.
\end{theorem}

\section{Proof intuition}

One implication is built into the definitions: the Poisson kernel is itself a
positive free pluriharmonic function. For the reverse implication, the key
input is that with only one genuinely free row, every positive free
pluriharmonic function is a completely positive image of the universal
Poisson kernel. A completely positive map preserves both inequalities in
(1), and it does not enlarge the constant. Thus the two sets of admissible
constants are identical, not merely comparable.

The one-row restriction enters only in the representation lemma below. Scalar
rows are commuting isometries and can be made unitary simultaneously by an
$\mathbb N^s$ direct limit. Fuglede's theorem then supplies the missing
adjoint commutation with the one nonscalar row.

\section{The representation input}

\begin{lemma}\label{lem:representation}
If at most one $n_i$ is larger than one, every positive free
$k$-pluriharmonic function $F$ on $\B_{\mathbf n}$, with coefficients in
$B(\E)$, has the form
\[
 F(\mathbf X)
 = (\mu\otimes\operatorname{id})
   [\boldsymbol{\Pcal}(\R,\mathbf X)]
\tag{2}
\]
for a completely positive map $\mu:C^*(\R)\to B(\E)$.
\end{lemma}

\begin{proof}
First normalize $F(0)=I_{\E}$. The positive-structure theorem of
\cite{pluri-source} gives $\K\supset\E$ and mutually commuting row
isometries $\mathbf V=(V_1,\ldots,V_k)$ whose compressed mixed moments are
the coefficients of $F$.

After relabeling, $V_1$ is the possible nonscalar row and
$V_2,\ldots,V_k$ are commuting scalar isometries $U_1,\ldots,U_s$. For
$p\le q$ in $\mathbb N^s$, connect a copy $\K_p$ of $\K$ to $\K_q$ by the
isometry $U^{q-p}$. On the Hilbert direct limit, define
\[
 \widetilde U_iJ_ph=J_pU_ih,
 \qquad
 \widetilde V_jJ_ph=J_pV_{1,j}h.
\]
These definitions are compatible with the connecting maps. Each
$\widetilde U_i$ is onto because
$\widetilde U_iJ_{p+e_i}h=J_ph$, hence is unitary. The
$\widetilde V_j$ remain isometries with orthogonal ranges. All extensions
commute; Fuglede's theorem \cite{conway} makes every $\widetilde V_j$
commute with every $\widetilde U_i^*$. Thus the extended rows doubly commute.

The universal property of the tensor product of the Cuntz--Toeplitz
algebras now gives a $*$-representation $\pi:C^*(\R)\to B(\widetilde\K)$
with $\pi(\R)$ equal to the extended tuple. Compression to $\E$ defines the
completely positive map $\mu(T)=J_{\E}^*\pi(T)J_{\E}$. Its moments are the
coefficients of $F$, so coefficient uniqueness yields (2).

For general $A=F(0)\ge0$, apply the normalized argument to
\[
 (A+\varepsilon I)^{-1/2}(F+\varepsilon I)
 (A+\varepsilon I)^{-1/2}.
\]
After conjugating back, the representing maps have norms
$\|A+\varepsilon I\|$. A point-ultraweak cluster limit as
$\varepsilon\downarrow0$ is completely positive and has the coefficients of
$F$, proving (2).
\end{proof}

\section{Proof of Theorem \ref{thm:metric}}

If $\mathbf A\sim_{H,c}\mathbf B$, apply the defining Harnack inequalities
to the positive free pluriharmonic function
$\mathbf X\mapsto\boldsymbol{\Pcal}(\R,\mathbf X)$. This gives (1), hence
$\mathbf A\sim_{\Pcal,c}\mathbf B$.

Conversely assume (1), and let $F$ be any positive free
$k$-pluriharmonic function with coefficients in $B(\E)$. By
Lemma~\ref{lem:representation}, choose a completely positive $\mu$ satisfying
(2). Apply the completely positive map
$\mu\otimes\operatorname{id}_{B(\mathcal H)}$ to all three terms of (1).
For every $0\le r<1$ this gives
\[
 c^{-2}F(r\mathbf B)\le F(r\mathbf A)\le c^2F(r\mathbf B).
\]
Thus $\mathbf A\sim_{H,c}\mathbf B$. The admissible constants are identical
in the two definitions, so their infima, their parts, and the resulting
logarithmic metrics are identical. \qed

\section{Verification, novelty, and limitations}

The proof is order-theoretic after Lemma~\ref{lem:representation}: no limiting
argument changes $c$. The representation lemma was independently packaged
and proof-checked earlier in this run; it is repeated here so that the present
packet is self-contained. The human reviewer should focus on the direct-limit
orientation, the use of Fuglede to obtain double commutation, and the tensor
order in (2).

A bounded search through 26 August 2026 used the exact source sentence, the
paper title and arXiv id, and combinations of ``Harnack'', ``Poisson'',
``regular polyball'', and ``one nonscalar row''. It found the source, later
applications, and a recent survey, but no resolution of the exact question or
this dimension class. The novelty verdict is therefore ``apparently new within
the bounded search''.

The theorem does not settle polyballs with two or more nonscalar rows. The
missing dilation would have to extend two commuting free row isometries while
preserving both orthogonal-range relations and cross-adjoint commutation;
ordinary $\mathbb N^s$ direct limits do not provide that extension.

\begin{thebibliography}{9}
\bibitem{metric-source}
G.~Popescu, \emph{Hyperbolic geometry on noncommutative polyballs},
J. Math. Anal. Appl. \textbf{456} (2017), 576--607;
arXiv:1701.00766.

\bibitem{pluri-source}
G.~Popescu, \emph{Free pluriharmonic functions on noncommutative polyballs},
Anal. PDE \textbf{9} (2016), 1185--1234;
arXiv:1511.08852.

\bibitem{conway}
J.~B. Conway, \emph{A Course in Functional Analysis}, second edition,
Graduate Texts in Mathematics 96, Springer, 1990. See the
Fuglede--Putnam theorem.
\end{thebibliography}
\end{document}
